% !TEX root = ../main.tex
\section{Apéndice: instrucciones de reproducibilidad}
\label{ap:reproducibilidad}

\subsection{Estructura del repositorio}

\begin{lstlisting}[caption={Disposición de archivos del proyecto.}]
mineria/
+- notebook.ipynb       Notebook principal (Jupyter / Google Colab)
+- README.md            Resumen del proyecto y guia de ejecucion
+- requirements.txt     Dependencias minimas con version
+- data/
|  +- penguins.csv      Dataset Palmer Penguins (CC-0)
+- images/              Figuras vectoriales generadas por el notebook
+- presentacion.pptx    Presentacion de 5 diapositivas
+- report/
   +- main.tex          Reporte LaTeX (este documento)
   +- sections/         Secciones modulares
   +- images/           Copia de figuras para el reporte
   +- references.bib    Bibliografia
\end{lstlisting}

\subsection{Ejecución local}

\begin{lstlisting}[language=bash,caption={Reproducción del experimento en local.}]
git clone <repo-url> mineria
cd mineria
python -m venv venv
venv\Scripts\activate              # Windows (en Linux/Mac: source venv/bin/activate)
pip install -r requirements.txt
jupyter notebook notebook.ipynb
\end{lstlisting}

\subsection{Ejecución en Google Colab}

\begin{enumerate}
  \item Subir \texttt{notebook.ipynb} a Google Drive.
  \item Abrirlo con Google Colab.
  \item Ejecutar \texttt{Entorno de ejecución $\rightarrow$ Ejecutar todo}.
  \item El dataset se descarga automáticamente desde el repositorio oficial si no está disponible localmente.
\end{enumerate}

\subsection{Compilación del reporte LaTeX}

\begin{lstlisting}[language=bash,caption={Compilación del PDF en el directorio report/.}]
cd report
pdflatex main.tex
bibtex main
pdflatex main.tex
pdflatex main.tex
\end{lstlisting}

\subsection{Checklist de Ciencia 2.0}

\begin{itemize}
  \item[\checkmark] Datos abiertos con licencia CC-0 y enlace al repositorio fuente.
  \item[\checkmark] Código abierto en notebook y scripts auxiliares.
  \item[\checkmark] Semillas aleatorias fijadas (\texttt{RANDOM\_STATE} = 42).
  \item[\checkmark] Dependencias declaradas con versiones mínimas (\texttt{requirements.txt}).
  \item[\checkmark] Notebook ejecutable de principio a fin sin pasos manuales.
  \item[\checkmark] Reporte LaTeX y presentación generados a partir de los mismos resultados.
\end{itemize}
